%% Question bank chapter file - generated by bank_tools.py build (do not edit while the bank is still being built:
%% it is regenerated from the per-test data in Question Bank\data). Every question is one bankq block:
%%   \begin{bankq}[id=..., source={...}, paper=.., qno=.., type=..., subtopic={...}, answer={...}, answer_src=key|solved,
%%                 working={...}, difficulty=1-10, diff_parts={K S I C F}, diff_note={...},
%%                 group=<passage id or empty>, also_in={duplicate copies in other tests}]
%% difficulty = min(10, K+S+I+C+F-1): K concepts 1-3, S steps 1-4, I insight 0-2, C computation 0-2, F format 0-1
%%   ...question LaTeX, options, figures (\bfig / \bfigi -> img/)...
%%   \end{bankq}
%% Passages shared by several questions are bankpassage blocks; their questions carry group=<passage id>.
\documentclass[10pt]{article}
\usepackage[a4paper,margin=1.8cm]{geometry}
\usepackage[T1]{fontenc}
\usepackage{amsmath,amssymb,mathtools}
\usepackage{graphicx}
\usepackage[export]{adjustbox}
\usepackage{xcolor,needspace,keyval}
\usepackage{enumitem}
\usepackage{booktabs,array,tabularx}
\usepackage[hidelinks]{hyperref}
\setlength{\parindent}{0pt}
\setlength{\parskip}{3pt}
\makeatletter
\define@key{bq}{id}{\def\bq@id{#1}}\define@key{bq}{source}{\def\bq@src{#1}}\define@key{bq}{paper}{\def\bq@paper{#1}}
\define@key{bq}{qno}{\def\bq@qno{#1}}\define@key{bq}{type}{\def\bq@type{#1}}\define@key{bq}{subtopic}{\def\bq@sub{#1}}
\define@key{bq}{answer}{\def\bq@ans{#1}}\define@key{bq}{answer_src}{\def\bq@asrc{#1}}\define@key{bq}{working}{\def\bq@work{#1}}
\define@key{bq}{group}{\def\bq@group{#1}}\define@key{bq}{also_in}{\def\bq@also{#1}}
\define@key{bq}{difficulty}{\def\bq@diff{#1}}\define@key{bq}{diff_parts}{\def\bq@dparts{#1}}\define@key{bq}{diff_note}{\def\bq@dnote{#1}}
\newenvironment{bankq}[1][]{\def\bq@work{}\def\bq@also{}\def\bq@group{}\def\bq@diff{?}\def\bq@dparts{}\def\bq@dnote{}\setkeys{bq}{#1}%
  \par\needspace{4\baselineskip}\medskip\noindent{\small\textbf{\bq@id}\hfill\color{gray}[\bq@src, Paper \bq@paper, Q\bq@qno]
  \bq@type\ $\cdot$ \bq@sub\ $\cdot$ \textbf{D\bq@diff}}\par\nopagebreak\relax}
  {\par{\small\color{teal}Answer: \textbf{\bq@ans}\ifx\bq@work\@empty\else\ (solved: \bq@work)\fi
  \ifx\bq@also\@empty\else\par\color{gray}Also in: \bq@also\fi}\par}
\define@key{bp}{id}{\def\bp@id{#1}}\define@key{bp}{source}{\def\bp@src{#1}}\define@key{bp}{questions}{\def\bp@qs{#1}}
\newenvironment{bankpassage}[1][]{\setkeys{bp}{#1}\par\needspace{4\baselineskip}\bigskip\noindent
  {\small\textbf{Passage \bp@id}\hfill\color{gray}[\bp@src; questions \bp@qs]}\par\nopagebreak\relax\itshape}{\par}
\makeatother
\newcommand{\bfig}[2]{\par\smallskip{\centering\includegraphics[width=#1pt,max width=0.62\linewidth,max height=0.32\textheight]{img/#2}\par}\smallskip}
\newcommand{\bfigi}[2]{\includegraphics[width=#1pt,max width=\linewidth,max height=0.32\textheight,valign=c]{img/#2}}
\newcommand{\bmissing}[1]{\textit{[figure missing in the source conversion: \detokenize{#1}]}}
\newenvironment{qoptions}{\begin{itemize}[leftmargin=1.8em,itemsep=1pt,topsep=2pt,label={}]}{\end{itemize}}
\newcommand{\opt}[2]{\item[(#1)] #2}
\providecommand{\fwstart}{}\providecommand{\fwend}{}\providecommand{\qfullsplit}{}
\newcolumntype{L}[1]{>{\raggedright\arraybackslash}p{#1}}
\newcolumntype{C}[1]{>{\centering\arraybackslash}p{#1}}
\newcolumntype{R}[1]{>{\raggedleft\arraybackslash}p{#1}}
\providecommand{\sectionhdr}[1]{}\providecommand{\subsectionhdr}[1]{}


\begin{document}
{\Large\bfseries EMI + AC --- Question Bank}\par\smallskip
{\small Physics chapter P12. Questions from all converted test papers, grouped by source series. Answers: printed key where it matches the paper, otherwise solved (answer\_src=solved, with working). Exact repeats across papers appear once, with also\_in listing the other copies.}\par\medskip\hrule

\section*{Score-I papers}
\begin{bankq}[id=SI12-P1-PHY-S1-Q02, source={Score-I - Test 12 (09 Mar 2022)}, paper=1, qno=2, type=multi, subtopic={LC oscillations}, answer={A,B,C}, answer_src=key, difficulty=2, diff_parts={K1 S1 I0 C0 F1}, diff_note={concepts: energy exchange in an LC circuit | steps: at maximum current the charge and electric field are zero and B is maximum; pick the incorrect statements}]
Consider the LC circuit below. At the time shown the current has its maximum value. At this time which of the following is/are incorrect?
\bfig{103.2}{SI12_p1_phy_s1_q02_1.png}
\begin{qoptions}
\opt{A}{The charge on the capacitor has its maximum value}
\opt{B}{The magnetic field inside the inductor is zero}
\opt{C}{The electric field inside the capacitor has its maximum value}
\opt{D}{The charge on the capacitor is zero}
\end{qoptions}
\end{bankq}
\begin{bankq}[id=SI12-P1-PHY-S3-Q10, source={Score-I - Test 12 (09 Mar 2022)}, paper=1, qno=10, type=numerical, subtopic={LC oscillations}, answer={6}, answer_src=key, difficulty=3, diff_parts={K1 S1 I1 C0 F1}, diff_note={concepts: energy conservation in an LC circuit | steps: omega = 1/sqrt(LC); Q^2 = q^2 + i^2/omega^2 gives n | I: use the energy relation}]
In an L-C circuit shown in the figure, $C=1$ F, $L=4$ H. At time $t=0$, charge in the capacitor is 4 C and it is decreasing at a rate of $\sqrt5$ C/s. The maximum charge in the capacitor is $n$ C, then value of $n$ will be:
\bfig{85.5}{SI12_p1_phy_s3_q10_1.png}
\end{bankq}
\begin{bankq}[id=SI12-P2-PHY-S2-Q06, source={Score-I - Test 12 (09 Mar 2022)}, paper=2, qno=6, type=numerical, subtopic={LCR circuits}, answer={20.00}, answer_src=key, difficulty=2, diff_parts={K1 S1 I0 C0 F1}, diff_note={concepts: AC and DC response of an LR circuit | steps: R from the DC current and Z from the AC current; XL = sqrt(Z^2 - R^2)}]
In the circuit shown, a 30 V d.c. source gives a current 2.0 A as recorded in the ammeter $A$ and 30 V a.c. source of frequency 100 Hz gives a current 1.2 A. The inductive reactance (in ohm) is :-
\bfig{108.9}{SI12_p2_phy_s2_q06_1.png}
\end{bankq}
\begin{bankq}[id=SI12-P2-PHY-S3-Q08, source={Score-I - Test 12 (09 Mar 2022)}, paper=2, qno=8, type=matching, subtopic={LCR circuits}, answer={P$\to$3, Q$\to$2, R$\to$1}, answer_src=key, difficulty=3, diff_parts={K1 S2 I0 C0 F1}, diff_note={concepts: phase relations in AC circuits | steps: phase of current in B relative to A; phase of voltage in D relative to B; phase of current in C relative to A}]
\bfig{153.9}{SI12_p2_phy_s3_q08_1.png}
\begin{center}\footnotesize
\begin{tabular}{|p{0.6cm}|p{4.8cm}|p{0.6cm}|p{1.6cm}|}
\hline
 & \centering List-I & \multicolumn{2}{c|}{List-II}\\
\hline
(P) & Phase of current in $B$ leads voltage of $A$ by a phase $\theta$. & (1) & $0<\theta<\dfrac\pi2$\\
\hline
(Q) & Phase of voltage in $D$ leads current of $B$ by phase $\theta$. & (2) & $-\dfrac\pi2<\theta<0$\\
\hline
(R) & Phase of current in $C$ leads current of $A$ by phase $\theta$. & (3) & $\theta=0$\\
\hline
\end{tabular}
\end{center}
\end{bankq}
\begin{bankq}[id=SI12-P2-PHY-S3-Q09, source={Score-I - Test 12 (09 Mar 2022)}, paper=2, qno=9, type=matching, subtopic={LCR circuits}, answer={P$\to$1, Q$\to$1, R$\to$3}, answer_src=key, difficulty=4, diff_parts={K2 S2 I0 C0 F1}, diff_note={concepts: induction in a coil circuit, discharging of a sphere | steps: branch x current for the coil circuit; potential difference for the sphere circuit; third circuit; match the properties}]
List-I contains some circuits and List-II contains some properties. In List-II current signifies instantaneous current and potential difference refers instantaneous potential difference.
\begin{center}\footnotesize
\begin{tabular}{|p{0.5cm}|p{4.3cm}|p{0.5cm}|p{1.5cm}|}
\hline
 & \centering List-I\newline Charge System & \multicolumn{2}{c|}{\shortstack{List-II\\Properties}}\\
\hline
(P) & \bfigi{89.4}{SI12_p2_phy_s3_q09_1.png}\rule{0pt}{32pt}\par A closely separated conducting wire is wrapped over an insulating cylinder. Wire is connected to a resistance $R$ and an alternating source via a switch $S$, which is closed at $t=0$ & (1) & In branch $x$, current varies\\
\hline
(Q) & \bfigi{115.5}{SI12_p2_phy_s3_q09_2.png}\rule{0pt}{27pt}\par A charged metallic sphere of radius $r$ is earthed with a resistance $R$ and a switch $S$, which is closed at $t=0$. & (2) & Between points $a$ and $b$ magnitude of potential difference increases always.\\
\hline
(R) & \bfigi{95.0}{SI12_p2_phy_s3_q09_3.png}\rule{0pt}{43pt}\par A circuit (used to measure an unknown resistance) is in balanced conditions (Galvanometer reads zero) & (3) & In branch $x$, current is constant\\
\hline
\end{tabular}
\end{center}
\end{bankq}
\end{document}
